\documentclass[10pt]{article}
\usepackage[T1,T2A]{fontenc}
\usepackage[utf8]{inputenc}
\usepackage[english,ukrainian]{babel}
\usepackage{geometry}
\usepackage{longtable}
\usepackage{array}
\usepackage{multirow}
\usepackage{hyperref}
\usepackage{mathptmx}
\renewcommand{\familydefault}{\rmdefault}

\geometry{a4paper,left=2cm,right=2cm,top=2cm,bottom=2cm}

\begin{document}

\begin{flushright}
ЗАТВЕРДЖЕНО\\
Наказ Державна Судова Адміністрація України\\
\_\_\_\_\_\_\_\_\_\_\_\_ 2026 року №\_\_\_\_
\end{flushright}

\begin{center}
\textbf{\large Галузевий профіль безпеки судової системи, що використовуються для надання хмарних послуг та/або послуг центру обробки даних публічним користувачам та/або критично важливим об’єктам інфраструктури}
\end{center}

\footnotesize
\begin{longtable}{|c|>{\raggedright\arraybackslash}p{4cm}|>{\raggedright\arraybackslash}p{6cm}|>{\raggedright\arraybackslash}p{2cm}|>{\raggedright\arraybackslash}p{3cm}|}
\hline
\textbf{№ з/п} & \textbf{Назва дії з безпеки інформації} & \textbf{Зміст дії} & \textbf{Заходи захисту} & \textbf{Мінімальні необхідні параметри}\\
\hline
\endfirsthead

\hline
\textbf{№ з/п} & \textbf{Назва дії з безпеки інформації} & \textbf{Зміст дії} & \textbf{Заходи захисту} & \textbf{Мінімальні необхідні параметри}\\
\hline
\endhead

\hline
\endfoot

\hline
\endlastfoot

\multicolumn{5}{|>{\raggedright\arraybackslash}p{16cm}|}{\textbf{УПРАВЛІННЯ ДОСТУПОМ (AC)}} \\ \hline
1 & НЕВДАЛІ СПРОБИ ВХОДУ В СИСТЕМУ - ОЧИЩЕННЯ АБО СТИРАННЯ МОБІЛЬНОГО ПРИСТРОЮ (AC-7(2)) & Інформація очищується або стирається з мобільних пристроїв на основі вимог або методів очищення або стирання після <AC-07(02)\_\allowbreak{}ODP[03 кількість> послідовних, невдалих спроб входу на пристрій & AC-7(2) & Параметр: ac\_\allowbreak{}7\_\allowbreak{}2\_\allowbreak{}01\newline Тип: integer\newline Значення: \textbf{3} \\ 
 &  & Визначається кількість послідовних невдалих спроб входу в систему до того, як інформація буде очищена або стерта з мобільних пристроїв &  & Параметр: ac\_\allowbreak{}7\_\allowbreak{}2\_\allowbreak{}odp\_\allowbreak{}03\newline Тип: integer\newline Значення: \textbf{3} \\ \hline
2 & АВТОМАТИЗОВАНЕ МАРКУВАННЯ (AC-15) & & AC-15 & \\ \hline
3 & АТРИБУТИ БЕЗПЕКИ ТА ПРИВАТНОСТІ (AC-16) & Визначено типи атрибутів безпеки, які мають бути пов'язані зі значеннями атрибутів безпеки для інформації, що зберігається, обробляється та/або передається; & AC-16 & Параметр: ac\_\allowbreak{}16\_\allowbreak{}c\_\allowbreak{}02\newline Тип: list\newline Значення: \textbf{default\_\allowbreak{}deny\_\allowbreak{}rule, abac\_\allowbreak{}rule\_\allowbreak{}1} \\ 
 &  & визначено значення атрибутів конфіденційності для типів атрибутів конфіденційності; &  &  \\ 
 &  & визначено системи, для яких мають бути встановлені дозволені атрибути безпеки; &  &  \\ 
 &  & визначено системи, для яких мають бути встановлені дозволені атрибути конфіденційності; &  &  \\ 
 &  & визначено атрибути безпеки, визначені як частина AC-16(a), які дозволені для систем; &  &  \\ 
 &  & визначено атрибути конфіденційності, визначені як частина AC-16(a), які дозволені для систем; &  &  \\ 
 &  & визначено значення атрибутів або діапазони для встановлених атрибутів; &  &  \\ 
 &  & визначено частоту, з якою слід переглядати атрибути безпеки на предмет відповідності; &  &  \\ 
 &  & визначено частоту, з якою слід переглядати атрибути конфіденційності на предмет відповідності; &  &  \\ 
 &  & надано засоби для асоціювання типів атрибутів безпеки зі значеннями атрибутів безпеки для інформації, що зберігається, обробляється та/або передається надано засоби для асоціювання типів атрибутів конфіденційності зі значеннями атрибутів конфіденційності для інформації, що зберігається, обробляється та/або передається виникають зв'язки з атрибутами; &  &  \\ 
 &  & зв'язки атрибутів зберігаються разом з інформацією; &  &  \\ 
 &  & на основі атрибутів, визначених у AC-16\_\allowbreak{}ODP[01] для систем, встановлюються наступні дозволені атрибути безпеки: атрибути безпеки; &  &  \\ 
 &  & на основі атрибутів, визначених у AC-16\_\allowbreak{}ODP[02] для систем, встановлюються наступні дозволені атрибути приватності: атрибути конфіденційності ; &  &  \\ 
 &  & AC-16(d) &  &  \\ 
 &  & Визначено наступні допустимі значення або діапазони атрибутів для кожного з встановлених атрибутів: значення або діапазони атрибутів; &  & Параметр: ac\_\allowbreak{}16\_\allowbreak{}f\_\allowbreak{}02\newline Тип: list\newline Значення: \textbf{default\_\allowbreak{}deny\_\allowbreak{}rule, abac\_\allowbreak{}rule\_\allowbreak{}1} \\ 
 &  & проводиться аудит змін до атрибутів; &  &  \\ 
 &  & атрибути безпеки перевіряються на відповідність частота; &  &  \\ 
 &  & атрибути конфіденційності перевіряються на відповідність частота &  &  \\ \hline
4 & КОНТРОЛЬ ДОСТУПУ ДЛЯ МОБІЛЬНИХ ПРИСТРОЇВ - ВИКОРИСТАННЯ ПИСЬМОВИХ ТА ПОРТАТИВНИЙ ПРИСТРОЇВ ДЛЯ ЗБЕРІГАННЯ ДАНИХ (AC-19(1)) & & AC-19(1) & \\ \cline{3-5}
 &  & КОНТРОЛЬ ДОСТУПУ ДЛЯ МОБІЛЬНИХ ПРИСТРОЇВ - ВИКОРИСТАННЯ ПЕРСОНАЛЬНИХ ПОРТАТИВНИХ ПРИСТРОЇВ ЗБЕРІГАННЯ ДАНИХ [Вилучено: Включено в MP-07]. AC-19(03) КОНТРОЛЬ ДОСТУПУ ДЛЯ МОБІЛЬНИХ ПРИСТРОЇВ - ВИКОРИСТАННЯ ПОРТАТИВНИХ ПРИСТРОЇВ ЗБЕРІГАННЯ ДАНИХ З НЕІДЕНТИФІКОВАНИМ ВЛАСНИКОМ [Вилучено: Включено в MP-07] & AC-19(2) & Параметр: ac\_\allowbreak{}19\_\allowbreak{}2\_\allowbreak{}01\newline Тип: list\newline Значення: \textbf{admin, security\_\allowbreak{}officer} \\ \cline{3-5}
 &  & Використання незахищених мобільних пристроїв на об'єктах, що містять системи, які обробляють, зберігають або передають секретну інформацію, заборонено, за винятком випадків, коли на це є спеціальний дозвіл уповноваженої посадової особи; & AC-19(4) & Параметр: ac\_\allowbreak{}19\_\allowbreak{}4\_\allowbreak{}a\newline Тип: list\newline Значення: \textbf{admin, security\_\allowbreak{}officer} \\ 
 &  & AC-19(04)(b)(01) заборона підключення незахищених мобільних пристроїв до систем з обмеженим доступом застосовується до осіб, яким уповноважена посадова особа дозволила використовувати незахищені мобільні пристрої на об'єктах, що містять системи, які обробляють, зберігають або передають інформацію з обмеженим доступом; &  &  \\ 
 &  & AC-19(04)(b)(02) дозвіл уповноваженої посадової особи на підключення незахищених мобільних пристроїв до незахищених систем вимагається від осіб, яким дозволено використовувати незахищені мобільні пристрої на об'єктах, що містять системи, які обробляють, зберігають або передають інформацію з обмеженим доступом; &  &  \\ 
 &  & AC-19(04)(b)(03) заборона використання внутрішніх або зовнішніх модемів чи бездротових інтерфейсів у складі незахищених мобільних пристроїв поширюється на осіб, яким уповноваженою посадовою особою дозволено використовувати незахищені мобільні пристрої під час виконання службових обов'язків, а також на осіб, які не мають права на використання таких пристроїв AC-19(04)(b)(04)[01] вибірковий огляд та перевірка незахищених мобільних пристроїв та інформації, що зберігається на них, посадовими особами є обов'язковими; &  &  \\ 
 &  & AC-19(04)(b)(04)[02] дотримання політики обробки інцидентів застосовується у разі виявлення секретної інформації під час випадкового огляду та перевірки незахищених мобільних пристроїв &  &  \\ 
 &  & Підключення захищенихмобільних пристроїв до засекречених систем обмежено відповідно до політики безпеки &  & Параметр: ac\_\allowbreak{}19\_\allowbreak{}4\_\allowbreak{}c\newline Тип: list\newline Значення: \textbf{default\_\allowbreak{}deny\_\allowbreak{}rule, abac\_\allowbreak{}rule\_\allowbreak{}1} \\ 
 &  & Визначено посадових осіб із захисту інформації, відповідальних за огляд та перевірку захищених мобільних пристроїв та інформації, що зберігається на цих пристроях &  & Параметр: ac\_\allowbreak{}19\_\allowbreak{}4\_\allowbreak{}odp\_\allowbreak{}01\newline Тип: list\newline Значення: \textbf{admin, security\_\allowbreak{}officer} \\ 
 &  & Визначено політики безпеки, що обмежують підключення захищених мобільних пристроїв до засекречених систем &  & Параметр: ac\_\allowbreak{}19\_\allowbreak{}4\_\allowbreak{}odp\_\allowbreak{}02\newline Тип: list\newline Значення: \textbf{default\_\allowbreak{}deny\_\allowbreak{}rule, abac\_\allowbreak{}rule\_\allowbreak{}1} \\ \hline
\multicolumn{5}{|>{\raggedright\arraybackslash}p{16cm}|}{\textbf{АУДИТ ТА ПІДЗВІТНІСТЬ (AU)}} \\ \hline
5 & НЕСПРОСТОВНІСТЬ (AU-10) & Надаються неспростовні докази того, що особа (або процес, що діє від імені особи) виконала дії & AU-10 & Параметр: au\_\allowbreak{}10\_\allowbreak{}01\newline Тип: list\newline Значення: \textbf{admin, security\_\allowbreak{}officer} \\ 
 &  & Визначено дії, на які поширюється принцип неспростовності &  & Параметр: au\_\allowbreak{}10\_\allowbreak{}odp\newline Тип: list\newline Значення: \textbf{login, logout, failed\_\allowbreak{}attempt} \\ \cline{3-5}
 &  & & AU-10(5) & \\ \hline
\multicolumn{5}{|>{\raggedright\arraybackslash}p{16cm}|}{\textbf{ІДЕНТИФІКАЦІЯ ТА АВТЕНТИФІКАЦІЯ (IA)}} \\ \hline
6 & ІДЕНТИФІКАЦІЯ ТА АВТЕНТИФІКАЦІЯ (КОРИСТУВАЧІВ ОРГАНІЗАЦІЇ) - БАГАТОФАКТОРНА АВТЕНТИФІКАЦІЯ ПРИВІЛЕЙОВАНИХ ОБЛІКОВИХ ЗАПИСІВ (IA-2(1)) & & IA-2(1) & \\ \cline{3-5}
 &  & & IA-2(2) & \\ \cline{3-5}
 &  & Приймаються та електронним шляхом підтверджуються повноваження облікових даних особистої ідентифікації & IA-2(12) & Параметр: ia\_\allowbreak{}2\_\allowbreak{}12\_\allowbreak{}01\newline Тип: list\newline Значення: \textbf{admin, security\_\allowbreak{}officer} \\ \hline
7 & УПРАВЛІННЯ АВТЕНТИФІКАТОРОМ - АВТЕНТИФІКАЦІЯ НА ОСНОВІ ВІДКРИТОГО КЛЮЧА (IA-5(2)) & & IA-5(2) & \\ \cline{3-5}
 &  & & IA-5(11) & \\ \cline{3-5}
 &  & Визначено вимоги до якості біометрії; & IA-5(12) & Параметр: ia\_\allowbreak{}5\_\allowbreak{}12\_\allowbreak{}odp\newline Тип: string\newline Значення: \textbf{автоматизований засіб моніторингу} \\ 
 &  & для біометричної автентифікації використовувати механізми, які задовольняють вимоги &  &  \\ \cline{3-5}
 &  & & IA-5(15) & \\ \cline{3-5}
 &  & Використовуються механізми виявлення атак із використанням штучно виготовлених артефактів для біометричних автентифікаторів & IA-5(17) & Параметр: ia\_\allowbreak{}5\_\allowbreak{}17\_\allowbreak{}01\newline Тип: string\newline Значення: \textbf{автоматизований засіб моніторингу} \\ \hline
8 & АВТЕНТИФІКАЦІЯ КРИПТОГРАФІЧНОГО МОДУЛЯ (IA-7) & Впроваджено механізми автентифікації в криптографічний модуль, який відповідає вимогам чинних законів, виконавчих розпоряджень, директив, політик, правил, стандартів та рекомендацій для такої автентифікації & IA-7 & Параметр: ia\_\allowbreak{}7\_\allowbreak{}01\newline Тип: list\newline Значення: \textbf{default\_\allowbreak{}deny\_\allowbreak{}rule, abac\_\allowbreak{}rule\_\allowbreak{}1} \\ \hline
\multicolumn{5}{|>{\raggedright\arraybackslash}p{16cm}|}{\textbf{ЗАХИСТ НОСІЇВ ІНФОРМАЦІЇ (MP)}} \\ \hline
9 & ЗНИЩЕННЯ ІНФОРМАЦІЇ НА НОСІЯХ ІНФОРМАЦІЇ (MP-6) & Застосовуються механізми очищення, надійність і цілісність яких відповідає категорії безпеки або рівню секретності інформації & MP-6 & Параметр: mp\_\allowbreak{}6\_\allowbreak{}b\newline Тип: string\newline Значення: \textbf{автоматизований засіб моніторингу} \\ \hline
10 & ЗНИЖЕННЯ КАТЕГОРІЇ БЕЗПЕКИ НОСІЇВ ІНФОРМАЦІЇ - ТАЄМНА ІНФОРМАЦІЯ (MP-8(4)) & Носії інформації, що містять інформацію з обмеженим доступом, знижуються в класі перед передачею особам, які не мають необхідних дозволів на доступ & MP-8(4) & Параметр: mp\_\allowbreak{}8\_\allowbreak{}4\_\allowbreak{}02\newline Тип: list\newline Значення: \textbf{admin, security\_\allowbreak{}officer} \\ \hline
\multicolumn{5}{|>{\raggedright\arraybackslash}p{16cm}|}{\textbf{ФІЗИЧНИЙ ЗАХИСТ ТА ЗАХИСТ НАВКОЛИШНЬОГО СЕРЕДОВИЩА (PE)}} \\ \hline
11 & МАРКУВАННЯ КОМПОНЕНТІВ (PE-22) & & PE-22 & \\ \hline
\multicolumn{5}{|>{\raggedright\arraybackslash}p{16cm}|}{\textbf{ЗАХИСТ СИСТЕМ ТА КОМУНІКАЦІЙ (SC)}} \\ \hline
12 & КОНФІДЕНЦІЙНІСТЬ ТА КРИПТОГРАФІЧНИЙ ЗАХИСТ (SC-8(1)) & & SC-8(1) & \\ \hline
13 & ВСТАНОВЛЕННЯ КЛЮЧАМИ (SC-12) & Встановлюються криптографічні ключі, коли в системі використовується криптографія відповідно до < SC-12\_\allowbreak{}ODP вимог > & SC-12 & Параметр: sc\_\allowbreak{}12\_\allowbreak{}01\newline Тип: string\newline Значення: \textbf{AES-256-GCM} \\ 
 &  & Здійснюється управління криптографічними ключами, коли в системі використовується криптографія, відповідно до вимог  &  & Параметр: sc\_\allowbreak{}12\_\allowbreak{}02\newline Тип: string\newline Значення: \textbf{AES-256-GCM} \\ \hline
14 & СЕРТИФІКАТИ ІНФРАСТРУКТУРИ ВІДКРИТИХ КЛЮЧІВ (SC-17) & & SC-17 & \\ \hline
15 & ЗАХИСТ ІНФОРМАЦІЇ В СТАНІ СПОКОЮ | КРИПТОГРАФІЧНИЙ ЗАХИСТ & Реалізовані криптографічні механізми для запобігання несанкціонованому розкриттю інформації, що знаходиться в стані спокою на системних компонентах або носіях & SC-28(1) & Параметр: sc\_\allowbreak{}28\_\allowbreak{}1\_\allowbreak{}01\newline Тип: string\newline Значення: \textbf{AES-256-GCM} \\ 
 &  & Реалізовані криптографічні механізми для запобігання несанкціонованій модифікації інформації, що знаходиться в стані спокою на системних компонентах або носіях &  & Параметр: sc\_\allowbreak{}28\_\allowbreak{}1\_\allowbreak{}02\newline Тип: string\newline Значення: \textbf{AES-256-GCM} \\ \hline
16 & ЕКРАНОВАНІ КАМЕРИ (SC-44) & & SC-44 & \\ \hline
17 & ПРИМУСОВЕ АПАРАТНЕ ЗАБЕЗПЕЧЕННЯ ВИКОНАННЯ (SC-49) & Впроваджено механізми апаратного розділення та застосування політик між доменами безпеки & SC-49 & Параметр: sc\_\allowbreak{}49\_\allowbreak{}01\newline Тип: list\newline Значення: \textbf{default\_\allowbreak{}deny\_\allowbreak{}rule, abac\_\allowbreak{}rule\_\allowbreak{}1} \\ 
 &  & Визначені домени безпеки, які потребують апаратного розділення та механізмів забезпечення дотримання політики &  & Параметр: sc\_\allowbreak{}49\_\allowbreak{}odp\newline Тип: list\newline Значення: \textbf{default\_\allowbreak{}deny\_\allowbreak{}rule, abac\_\allowbreak{}rule\_\allowbreak{}1} \\ \hline
18 & АПАРАТНИЙ ЗАХИСТ (SC-51) & Визначені уповноважені особи, які повинні виконувати процедури вимкнення та повторного ввімкнення апаратного захисту від запису; & SC-51 & Параметр: sc\_\allowbreak{}51\_\allowbreak{}odp\_\allowbreak{}02\newline Тип: list\newline Значення: \textbf{admin, security\_\allowbreak{}officer} \\ 
 &  & SC-51a. використовується апаратний захист від запису для компонентів мікропрограми системи; &  &  \\ 
 &  & SC-51b.[01] впроваджено спеціальні процедури для уповноважених осіб для ручного вимкнення апаратного захисту від запису для модифікацій мікропрограми; &  &  \\ 
 &  & SC-51b.[02] реалізовано спеціальні процедури для уповноважених осіб для повторного увімкнення захисту від запису перед поверненням до робочого режиму &  &  \\ \hline
\multicolumn{5}{|>{\raggedright\arraybackslash}p{16cm}|}{\textbf{ЦІЛІСНІСТЬ СИСТЕМИ ТА ІНФОРМАЦІЇ (SI)}} \\ \hline
19 & ЦІЛІСНІСТЬ ПРОГРАМНОГО ЗАБЕЗПЕЧЕННЯ, ПРОГРАМНОГО ЗАБЕЗПЕЧЕННЯ ТА ІНФОРМАЦІЇ (SI-7) & Визначені дії, яких слід вжити при виявленні несанкціонованих змін у програмному забезпеченні & SI-7 & Параметр: si\_\allowbreak{}7\_\allowbreak{}odp\_\allowbreak{}04\newline Тип: list\newline Значення: \textbf{login, logout, failed\_\allowbreak{}attempt} \\ 
 &  & Визначені дії, яких слід вжити несанкціонованих змін у прошивці; &  & Параметр: si\_\allowbreak{}7\_\allowbreak{}odp\_\allowbreak{}05\newline Тип: list\newline Значення: \textbf{login, logout, failed\_\allowbreak{}attempt} \\ 
 &  & при виявленні SI-07\_\allowbreak{}ODP[06] визначені дії, яких слід вжити несанкціонованих змін до інформації; &  &  \\ 
 &  & при виявленні SI-07a.[01] використовуються засоби перевірки цілісності для виявлення несанкціонованих змін у програмному забезпеченні; &  &  \\ 
 &  & SI-07a.[02] використовуються засоби перевірки цілісності для виявлення несанкціонованих змін у мікропрограмі; &  &  \\ 
 &  & SI-07a.[03] використовуються засоби перевірки цілісності для виявлення несанкціонованих змін до інформації; &  &  \\ 
 &  & SI-07b.[01] виконуються дії при виявленні несанкціонованих змін у програмному забезпеченні; &  &  \\ 
 &  & SI-07b.[02] виконуються дії несанкціонованих змін у прошивці; &  &  \\ 
 &  & при виявленні SI-07b.[03] виконуються дії несанкціонованих змін в інформації. при виявленні &  &  \\ \hline


\end{longtable}

\end{document}
